% {{FOLDER}}/main.tex — {{TITLE}}
% texref · paper template (two columns) · created {{CREATED}}
% Build: make · make launch (Skim watch) · make name NAME=<pdf> · make archive TAG=<tag>
% Two-column notes: xltabular/longtable cannot be used; use tabularx on \columnwidth inside table,
% or table*/figure* (with \textwidth) to span both columns. See {{UNIT_DIR}}/{{FIRST_UNIT_FILE}}.
% Default LaTeX fonts, as in the other house templates. Style macros (defined in the paperstyle block): \deck{..} \lead{X}{..} \panelhead{..} \readingguide{..} \eqnote{..}
%   \pill{A} \pillaccent{B} \colophon{..}   environments: summary · eqpanel · eqhero · figframe · pullquote · papertable
% =====================================================================================================================
\documentclass[10pt,letterpaper,twocolumn]{article}
\usepackage[margin=0.8in,columnsep=24pt]{geometry}
% ---------------------------------------------------------------- core packages
\usepackage{amsmath,amssymb,amsthm,mathtools}
\usepackage{microtype}
\usepackage[dvipsnames,table]{xcolor}
\usepackage{graphicx}
\graphicspath{{figures/}}
\usepackage{setspace}
\usepackage{titlesec}
\usepackage{enumitem}
\setlist{itemsep=3pt,topsep=4pt,leftmargin=1.35em}
\usepackage{parskip}
\usepackage{caption}
\usepackage{listings}
\usepackage{needspace}
{{INCLUDE:tables}}
{{INCLUDE:code}}
{{INCLUDE:paperstyle}}
\usepackage{bm}
\setstretch{1.08}
\setlength{\columnseprule}{0.3pt}
\makeatletter   % color the column rule (the kernel draws it with \normalcolor; multicol's \columnseprulecolor does not apply)
\patchcmd{\@outputdblcol}{\normalcolor\vrule}{\color{hairtwo}\vrule}{}{}
\makeatother

% ---------------------------------------------------------------- document metadata (edit here)
\newcommand{\doctitle}{{{TITLE}}}
\newcommand{\docshorttitle}{\doctitle}          % running header; shorten if the title is long
\newcommand{\docsubtitle}{{{SUBTITLE}}}
\newcommand{\doceyebrow}{{{TAGLINE}}}           % small letterspaced line above the title, e.g. "Frontier · Diligence Brief"
\newcommand{\docauthor}{{{AUTHOR}}}
\newcommand{\docdate}{{{DATE}}}
\newcommand{\docversion}{{{VERSION}}}           % prints faded at the foot of every page

% ---------------------------------------------------------------- theorems
\newtheorem{definition}{Definition}[section]
\newtheorem{proposition}{Proposition}[section]
\newtheorem{assumption}{Assumption}[section]

% Title block + abstract spanning both columns
\makeatletter
\newcommand{\spanningtop}[1]{\twocolumn[\begin{@twocolumnfalse}#1\end{@twocolumnfalse}]}
\makeatother
% Sections flow in the columns. Put \unitpage before a \section to start it on a fresh page (rule suppressed).
\newcommand{\unitpage}{\clearpage\secrulefalse}

{{INCLUDE:macros}}
{{INCLUDE:bib}}
\renewcommand*{\bibfont}{\small}
{{INCLUDE:hyperref}}
\hypersetup{linkcolor=accent,citecolor=accent,urlcolor=accent}

\begin{document}

% =====================================================================================================================
\spanningtop{%
\begin{center}
\ifdefempty{\doceyebrow}{}{{\sffamily\scriptsize\bfseries\color{ink3}\textls[180]{\MakeUppercase{\doceyebrow}}\par}\vspace{14pt}}
{\LARGE\bfseries \doctitle\par}
\vspace{6pt}
{\large\itshape\color{ink2} \docsubtitle\par}
\vspace{10pt}
{\normalsize \docauthor\par}
\vspace{3pt}
{\small\color{ink3} \docdate\par}
\end{center}
\vspace{10pt}
\input{{{UNIT_DIR}}/abstract}
\vspace{10pt}
}
\thispagestyle{plain}
\secrulefalse   % the first section sits under the abstract panel without a hairline

% =====================================================================================================================
% >>> texref:units  (managed by tex-add / tex-rm — edit order here if you must, keep the markers)
{{UNIT_INPUTS}}
% <<< texref:units

% =====================================================================================================================
\startappendices
% >>> texref:appendices
{{APPENDIX_INPUTS}}
% <<< texref:appendices

% =====================================================================================================================
\printbibliography

\colophon{\textbf{Colophon.} \doctitle, version \docversion.}

\end{document}
